\documentclass[10pt]{amsart}
\usepackage[a4paper,margin=22mm]{geometry}
\usepackage[no-math]{fontspec}
\setmainfont{Latin Modern Roman}[SmallCapsFont={lmromancaps10-regular.otf}]
\usepackage{amsmath,amssymb,amsthm,mathtools,hyperref}
\usepackage[shortlabels]{enumitem}
\emergencystretch=3em
\newtheorem{theo}{Theorem}[section]
\newtheorem{lemm}[theo]{Lemma}
\newtheorem{prop}[theo]{Proposition}
\newtheorem{coro}[theo]{Corollary}
\theoremstyle{definition}
\newtheorem{defi}[theo]{Definition}
\newtheorem{rema}[theo]{Remark}
\newtheorem{exam}[theo]{Example}
\providecommand{\ie}{i.e.\ }
\providecommand{\eg}{e.g.\ }
\newcommand{\PP}{\mathbb{P}}
\newcommand{\N}{\mathbb{N}}
\newcommand{\Z}{\mathbb{Z}}
\newcommand{\R}{\mathbb{R}}
\newcommand{\Cbb}{\mathbb{C}}

\newcommand{\V}{\mathcal{V}}
\newcommand{\I}{\mathcal{I}}
\newcommand{\J}{\mathcal{J}}
\newcommand{\cL}{\mathcal{L}}
\newcommand{\CC}{\mathcal{C}}
\newcommand{\B}{\mathcal{B}}
\newcommand{\cP}{\mathcal{P}}

\newcommand{\x}{\mathbf{x}}
\newcommand{\p}{\mathbf{p}}
\newcommand{\ol}{\overline}
\newcommand{\ii}{\textnormal{i}}

\DeclareMathOperator{\trop}{trop}
\DeclareMathOperator{\In}{in}
\DeclareMathOperator{\inv}{inv}
\DeclareMathOperator{\supp}{supp}
\DeclareMathOperator{\spanrm}{span}
\DeclareMathOperator{\Zar}{Zar}
\DeclareMathOperator{\rank}{rank}
\DeclareMathOperator{\cone}{cone}
\DeclareMathOperator{\facets}{facets}
\DeclareMathOperator{\dd}{d}
\DeclareMathOperator{\link}{link}
\DeclareMathOperator{\argmin}{argmin}
\DeclareMathOperator{\argmax}{argmax}
\DeclareMathOperator{\rowspan}{rowspan}
\DeclareMathOperator{\Imrm}{Im}
\DeclareMathOperator{\rec}{rec}



\newcommand{\reldp}{\mathrel{:}}

\def\dashmapsto{\mapstochar\dashrightarrow}

%%%%%%%%%%%%%%%%%%%%%%%%%%%%%%%%%%%%%
\newcommand*{\mk}{\mkern -1mu}
\newcommand*{\Mk}{\mkern -2mu}
\newcommand*{\mK}{\mkern 1mu}
\newcommand*{\MK}{\mkern 2mu}

\hypersetup{urlcolor=purple, linkcolor=blue, citecolor=red}


\newcommand*{\romanenumi}{\renewcommand*{\theenumi}{\roman{enumi}}}
\newcommand*{\Romanenumi}{\renewcommand*{\theenumi}{\Roman{enumi}}}
\newcommand*{\alphenumi}{\renewcommand*{\theenumi}{\alph{enumi}}}
\newcommand*{\Alphenumi}{\renewcommand*{\theenumi}{\Alph{enumi}}}
\title[Universal circuit bases for partial inversion]{Universal circuit Gr\"obner bases for partial coordinate inversions of complex linear spaces}
\author{Georgy Scholten}
\author{Cynthia Vinzant}
\date{Algebraic Combinatorics2(4)(2019),645–661; DOI10.5802/alco.65}
\hypersetup{hidelinks}
\begin{document}
\maketitle
\section{Introduction}
In this paper we extend the results of Proudfoot and Speyer to the image of a linear space $\cL\subset \Cbb^n$ under inversion of some subset of coordinates.
For $I\subseteq \{1,\ldots , n\}$, consider the rational map $\inv_I: \Cbb^n \dashrightarrow \Cbb^n $ defined by 
\[(\inv_I(x))_i =
\begin{cases}
1/x_i &\text{ if } i\in I\\
x_i &\text{ if } i\not\in I.\\
\end{cases}
\]
Let $\inv_I(\cL)$ denote the Zariski-closure of the image of $\cL$ under this map, which is an affine variety in $\Cbb^n$. One can interpret $\inv_I(\cL)$ as 
an affine chart of the closure of $\cL$ in the product of projective spaces $(\PP^1)^n$, as studied in \cite{AB16}, or as the projection 
of the graph of $\cL$ under the map $x\dashmapsto \inv_{[n]}(x)$, studied in \cite{FSW17}, onto complementary subsets of the $2n$ coordinates. 
We give a degeneration of the coordinate ring of $\inv_I(\cL)$
to the Stanley--Reisner ring of a simplicial complex generalizing the broken circuit complex of a matroid. 
This involves constructing a universal Gr\"obner basis for the ideal of polynomials vanishing on $\inv_I(\cL)$.

Let $\Cbb[\x]$ denote the polynomial ring $\Cbb[x_1, \ldots ,x_n]$ and for any $\alpha \in (\Z_{\geq 0})^n$, let 
$\x^{\alpha}$ denote $\prod_{i=1}^n x_i^{\alpha_i}$. For a subset $S\subseteq[n]$, we will also use $\x^S$ to denote $\prod_{i\in S}x_i$. 
As in \cite{PS}, the \emph{circuits} of the matroid $M(\cL)$ corresponding to $\cL$ give rise to a universal Gr\"obner basis for the 
ideal of polynomials vanishing on $\inv_I(\cL)$. 
We say that a linear form $\ell(x)= \sum_{i\in [n]} a_i x_i$ vanishes on $\cL$ if $\ell(x) = 0$ for all $x\in \cL$. 
The support of $\ell$, $\supp(\ell)$, is $\{i\in [n] : a_i\neq 0\}$. The minimal supports of nonzero linear forms vanishing on $\cL$
are called \emph{circuits} of the matroid $M(\cL)$ and for every circuit $C\subset [n]$, there is a unique (up to scaling) 
linear form $\ell_C= \sum_{i\in C} a_i x_i$ vanishing on $\cL$ with support $C$. To each circuit, we associate the polynomial 
\begin{equation}\label{eq:fC_def}
f_C(\x) = \x^{C\cap I} \cdot \ell_C(\inv_I(\x)) = \sum_{i\in C\cap I} a_i \x^{C\cap I\backslash\{ i\}} + \sum_{i\in C\backslash I } a_i \x^{C\cap I \cup\{ i\}}.
\end{equation}


\begin{theo}\label{thm:main}
Let $\cL\subseteq \Cbb^n$ be a $d$-dimensional linear space and 
let $\I \subseteq \Cbb[\x]$ be the ideal of polynomials vanishing on $\inv_I(\cL)$. 
Then $\{f_C : C \text{ is a circuit of }M(\cL)\}$ is a universal Gr\"obner basis for $\I$. 
For $w\in (\R_+)^n$ with distinct coordinates, the initial ideal $\In_w(\I)$ is the Stanley--Reisner 
ideal of the semi-broken circuit complex $\Delta_w(M(\cL), I)$. 
\end{theo}
The simplicial complex $\Delta_w(M(\cL), I)$ will be defined in Section~\ref{sec:SimplicialComplex}.
\section{Background}\label{Background}

In this section, we review the necessary background on Gr\"obner bases, simplicial complexes, Stanley--Reisner ideals, matroids, and 
previous research on reciprocal linear spaces. 

\subsection{Gr\"obner bases and degenerations} \label{subsec:Grobner}

A finite subset $F$ of an ideal $\I \subset \Cbb[\x]$ is a \emph{universal Gr\"obner basis} for $\I$ 
if it is a Gr\"obner basis with respect to every monomial order on $\Cbb[\x]$. 
An equivalent definition using weight vectors is given as follows. For $w\in (\R_{\geq 0})^n$ and $f = \sum_\alpha c_{\alpha} \x^{\alpha} \in \Cbb[\x]$, 
define the degree and initial form of $f$ with respect to $w$ to be 
\[
\deg_w(f) = \max\{w^T\alpha : c_{\alpha}\neq 0\} \quad \text{and} \quad
\In_w(f) = \sum_{\alpha: w^T\alpha = \deg_w(f) } c_{\alpha} \x^{\alpha}. \]
The initial ideal $\In_w(\I)$ of an ideal $\I$ is the ideal generated by initial forms of polynomials in $\I$, \ie $\In_w(\I) = \langle \In_w(f) : f\in \I\rangle$. 
Then $F\subset \I$ is a universal Gr\"obner basis for $\I$ if and only if for every $w\in (\R_{\geq 0})^n$, the polynomials $\In_w(F)$ generate $\In_w(\I)$. 
See \cite[Chapter~1]{GrobnerPolytope}. 

For homogeneous $\I$, $\In_w(\I)$ is a \emph{flat degeneration} of $\I$. 
For $f\in \Cbb[\x]$ and an integer vector $w\in (\Z_{\geq 0})^n$, define 
\[t^w\cdot f = t^{\deg_w(f)}f(t^{-w_1}x_1, \ldots , t^{-w_n}x_n) \in \Cbb[t,\x]
\quad \text{and} \quad
\ol{\I}^w = \langle t^{w}\cdot f : f\in \I \rangle \subset \Cbb[t,\x].
\]
The ideal $\ol{\I}^w$ defines a variety in $\mathbb{A}^1(\Cbb) \times \PP^{n-1}(\Cbb)$,
namely the Zariski-closure 
\[
\V(\ol{\I}^w) = \ol{\bigl\{ (t, [t^{w_1} x_1 : \ldots : t^{w_n}x_n ]) \ \text{ such that } \ t\in {\Cbb^*}, x\in \V(\I)\bigl\}}^{\Zar}.
\]
Letting $t$ vary from $1$ to $0$ gives a flat deformation from $\V(\I)$ to 
the variety $\V(\In_w(\I))$. Formally, for any $\gamma\in \Cbb$, let $\ol{\I}^w(\gamma)$ denote the ideal in $\Cbb[\x]$
obtained by substituting $t = \gamma$. 
Then $\ol{\I}^w(1)$ equals $ \I$, $\ol{\I}^w(0)$ equals $\In_w(\I)$, 
and for $\gamma\in \Cbb^*$, the variety of $\ol{\I}^w(\gamma)$ 
consists of the points $\{[\gamma^{w_1} x_1 : \ldots : \gamma^{w_n}x_n] : x\in \V(\I)\}$. 
We note that the ideal $\ol{\I}^w(\gamma)$ is well-defined for any $w\in \R^n$. 
All the ideals $\ol{\I}^w(\gamma)$ have the same Hilbert series. 
In particular, taking $\gamma = 0,1$ shows that $\I$ and $\In_w(\I)$ have the same Hilbert series. 


\subsection{Simplicial complexes and Stanley--Reisner ideals} \label{subsec:SRI}
\looseness-1
A Stanley--Reisner ideal is a square-free monomial ideal. Its combinatorial 
properties are governed by a simplicial complex. 
A simplicial complex $\Delta$ on vertices $\{1,\ldots ,n\}$ is a collection of subsets of 
$\{1,\ldots , n\}$, called faces, that is closed under taking subsets. If $S \in \Delta$ has 
cardinality $k+1$, we call it a face of dimension $k$. A \emph{facet} of $\Delta$ is a face maximal in $\Delta$ under inclusion. 
Given a simplicial complex $\Delta$ on $[n]\backslash \{i\}$, define the cone of $\Delta$ over $i$ to be 
\[ 
\cone (\Delta, i) = \Delta \cup \{ S\cup \{i\} : S \in \Delta\}, \]
which is a simplicial complex on $[n]$ whose facets are in bijection with the facets of $\Delta$. 

\begin{defi}[See \eg {\cite[Chapter~II]{StanleyCCABook}}]Let $\Delta$ be a simplicial complex on vertices $\{1,\ldots , n\}$. The \emph{Stanley--Reisner ideal} of $\Delta$ is the square-free monomial ideal 
\[
\I_{\Delta} = \left\langle \x^S \colon S\subseteq [n], \ S\not\in \Delta \right\rangle
\]
generated by monomials corresponding to the non-faces of $\Delta$. The \emph{Stanley--Reisner ring} of $\Delta$ is the quotient ring $\Cbb[\x]/\I_{\Delta}$. 
\end{defi}

The ideal $\I_{\Delta}$ is radical 
and it equals the intersection of prime ideals
\[
\I_{\Delta} = \bigcap_{F \text{ a facet of }\Delta} \langle x_i : i\not\in F \rangle.
\]
This writes the variety $\V(\I_{\Delta})$ as the union of coordinate subspaces $\spanrm \{e_i : i\in F\}$ 
where $F$ is a facet of $\Delta$. In particular, if $\Delta$ has $k$ facets of dimension $d-1$, then 
$\V(\I_{\Delta})\subseteq \PP^{n-1}(\Cbb)$ is a variety of dimension $d-1$ and degree $k$. See \cite[Chapter~II]{StanleyCCABook}. 


\subsection{Matroids}
Matroids are a combinatorial model for many types of independence relations. See \cite{Oxley} for 
general background on matroid theory. We can associate a matroid $M(\cL)$ to a linear space $\cL\subset \Cbb^n$ as follows.
Write a $d$-dimensional linear space $\cL\subset \Cbb^n$ as the rowspan of a $d\times n$ matrix $A = (a_1, \ldots , a_n)$. 
A set $I \subseteq [n]$ is \emph{independent} in $M(\cL)$ 
if the vectors $\{a_i : i\in I\}$ are linearly independent in $\Cbb^d$. For any invertible matrix $U\in \Cbb^{d\times d}$, the vectors $\{a_i : i\in I\}$
are linearly independent if and only if the vectors $\{Ua_i : i\in I\}$ are also independent, implying that this condition is independent of the choice of basis for $\cL$. 
Indeed, $I\subseteq [n]$ is independent in $M(\cL)$ if and only if the coordinate linear forms $\{x_i : i\in I\}$ are linearly independent when restricted to $\cL$. 


Maximal independent sets are called \emph{bases} and minimal dependent sets are called \emph{circuits}. We use $\B(M)$ and $\CC (M)$ to denote the 
set of bases and circuits of a matroid $M$, respectively. An element $i\in [n]$ is called a \emph{loop} if $\{i\}$ is a circuit, and a \emph{co-loop} if~$i$ is contained in 
every basis of $M$. 
The \emph{rank} of a subset $S \subseteq [n]$ is the largest size of an independent set in $S$. A \emph{flat} is a set $F\subseteq [n]$ that is maximal for its rank, meaning that 
$\rank(F) <\rank (F \cup \{i\}) $ for any $i\not\in F$. 


Let $M$ be a matroid on $[n]$ and $i \in [n]$. The \emph{deletion} of $M$ by $i$, denoted $M\backslash i $, is the matroid
on the ground set $[n]\backslash i$ whose independent sets are subsets $I\subset [n]\backslash i$ that are independent in $M$. 
If $i$ is not a co-loop of $M$, then 
\[
\B(M\backslash i) = \{ B\in \B (M): i\notin B \} 
\quad \text{and} \quad
\CC (M\backslash i) = \{C\in \CC (M):i\notin C \}. 
\]
More generally, the deletion of $M$ by a subset $S\subset [n]$, denoted $M\backslash S$, is the matroid obtained from $M$ by 
successive deletion of the elements of $S$. The \emph{restriction} of $M$ to a subset $S$, denoted $M|_S$, is 
the deletion of $M$ by $[n]\backslash S$. 


If $i$ is not a loop of $M$, then the \emph{contraction} of $M$ by $i$, denoted $M/ i $, is the matroid
on the ground set $[n]\backslash \{i\}$ whose independent sets are subsets $I\subset [n]\backslash i$ 
for which $I\cup \{i\}$ is independent in $M$. 
Then 
\begin{align*}
\B(M/ i) &= \{B\backslash i \reldp B \in \B (M), i\in B \}, \text{ and } \\
\CC (M/ i) & = \text{inclusion minimal elements of }\{C\backslash i \reldp C\in \CC (M) \}. 
\end{align*}
If $i$ is a loop of $M$, then we define the contraction of $M/i$ to be the deletion $M\backslash i$. 
The contraction of $M$ by a subset $S\subset [n]$, denoted $M/ S$, is obtained from $M$ by 
successive contractions by the elements of $S$. 


For linear matroids, deletion and contraction correspond to projection and intersection in the following sense. 
For $S\subset [n]$, let $\cL\backslash S$ denote the linear subspace of $\Cbb^{[n]\backslash S}$ obtained by projecting $\cL$ away 
from the coordinate space $\Cbb^S = \spanrm \{e_i : i\in S\}$. 
Let $\cL/S$ denote the intersection of $\cL$ with $\Cbb^{[n]\backslash S}$. Then 
\[ M(\cL)\backslash S = M(\cL\backslash S) 
\quad \text{and} \quad 
M(\cL)/S = M(\cL/S). 
\]

Many interesting combinatorial properties of a matroid can be extracted from a simplicial complex 
called the broken circuit complex. Given a matroid $M$ and the usual ordering $1<2< \cdots <n$ on $[n]$, 
a \emph{broken-circuit} of $M$ is a subset of the form $C\backslash \min(C)$ where $C\in \CC (M)$. 
The \emph{broken-circuit complex} of $M$ is the simplicial complex on $[n]$ whose faces are 
the subsets of $[n]$ not containing any broken circuit. 

\section{A semi-broken circuit complex}\label{sec:SimplicialComplex}

Let $M$ be a matroid on elements $[n]$ and suppose $I \subseteq [n]$. 
A vector $w \in \R^n$ with distinct coordinates gives an ordering on $[n]$, where $i<j$ whenever $w_i< w_j$. 
Without loss of generality, we can assume $w_1< \cdots <w_n$, which induces the usual order $1<\cdots <n$. 
Given a circuit $C$ of $M$ we define an \emph{$I$-broken circuit} of $M$ to be 
\[
b_I(C)  =  
\begin{cases}
C\backslash\min(C) & \text{if }C \subseteq I \\
(C\cap I) \cup \max(C\backslash I) & \text{if }C \not\subseteq I.
\end{cases}
\]
Now we define the \emph{$I$-broken circuit complex} of $M$ to be 
\begin{equation}\label{eq:complexDef}
\Delta_w(M,I)  =  \{ \tau \subseteq [n] \reldp \tau \text{ does not contain an $I$-broken circuit of $M$}\}.
\end{equation}
Note that an $[n]$-broken circuit is a broken circuit in the usual sense and 
$\Delta_w(M,[n])$ is the well-studied broken circuit complex of $M$. 



\begin{exam} \label{ex:NonGenMatroid} 
Consider the rank-$3$ matroid on $[5]$ with circuits $\CC = \{124, 135, 2345\}$.
Let $I = \{1,2,3\}$ and suppose $w\in (\R_+)^5$ with $w_1<\cdots < w_5$.
Then its $I$-broken circuits are 
$b_I(124) = 124$, $b_I(135) = 135$, and
$b_I(2345) = 235$.
The simplicial complex $\Delta_w(M,I)$ is a pure $2$-dimensional simplicial complex with facets:
\[ \facets(\Delta_w(M,I))  =  \{ 123, 125, 134, 145, 234, 245, 345 \}.\]
\end{exam}


The $I$-broken circuit complex shares many properties with the classical one, which will imply that 
it is always pure of dimension $\rank (M)-1$. 


\begin{theo}\label{thm:Complex}
Let $\Delta_w(M,I)$ be the $I$-broken circuit complex defined in \eqref{eq:complexDef}. 
\begin{enumerate}[(\alph*)]
\item\label{theo3.2_a} %[(a)] 
If $i\in I$ is a loop of $M$, then $\Delta_w(M,I) = \emptyset$. 
\item\label{theo3.2_b} %[(b)] 
If $i\in I$ is a coloop of $M$, then $\Delta_w(M,I)=\cone (\Delta_w(M/i,I\backslash i),i)$.
\item\label{theo3.2_c} %[(c)] 
If $i=\max(I)$ is neither a loop nor a coloop of $M$, then 
\[\Delta_w(M,I) = \Delta_w(M\backslash i,I \backslash i)  \cup  \cone ( \Delta_w(M/i,I\backslash i), i). \]
\end{enumerate}
\end{theo}

\begin{proof}%\hskip0pt

%\begin{enumerate}[(\alph*)]
%\item 
\ref{theo3.2_a}~%
If $i\in I$ is a loop, then $C = \{i\}$ is a circuit of $M$ with $b_I(C) = \emptyset$. 

%\item %
\ref{theo3.2_b}~%
If $i\in I$ is a coloop, then no circuit of $M$, and hence no $I$-broken circuit, contains $i$. 
The circuits of $M$ are exactly the circuits of the contraction $M/i$ and the $I$-broken circuits of $M$ 
are the $(I\backslash i)$-broken circuits of $M/i$. 
Therefore $\tau$ is a face of $\Delta_w(M,I)$ if any only if $\tau\backslash i$ is a face of $\Delta_w(M/i,I\backslash i)$.

%\item 
\ref{theo3.2_c}~%
($\subseteq$) Let $\tau$ be a face of $\Delta_w(M,I)$. We will show that if $i\not\in \tau$, then $\tau$ is a 
face of $\Delta_w(M\backslash i ,I\backslash i)$ and if $i \in \tau$, then $\tau\backslash i$ is a face of $\Delta_w(M/i ,I\backslash i)$.

If $i\not\in \tau$ and $C$ is a circuit of the deletion $M\backslash i$, then $C$
is a circuit of $M$, and $b_I(C)=b_{I\backslash i}(C)$ is an $I$-broken circuit of $M$ and therefore is not contained in $\tau$. 
If $i \in \tau$ and $C$ is a circuit of the contraction $M/i$, then either $C$ or $C\cup \{i\}$ is a circuit of $M$. 
In the first case, we again have that $b_I(C)=b_{I\backslash i}(C)$ is not contained in $\tau$ and thus not contained in $\tau\backslash i$. 
Secondly, suppose that $C\cup\{i\}$ is a circuit of $M$. 
If $C\subseteq I$, then $b_I(C\cup\{i\})$ is equal to $C\cup\{i\}\backslash \min(C\cup\{i\})$. Since $i$ is the maximum element of $I$, 
this equals $C\backslash \min(C) \cup\{i\}$. This set is not contained in $\tau$. Therefore $b_{I\backslash i}(C) = C\backslash \min(C)$ 
is not contained in $\tau \backslash i$. 
If $C\not\subseteq I$, then the $I$-broken circuit of $C\cup \{i\}$ is $(C\cap I) \cup \{ i\} \cup \max(C\backslash I)$, 
which equals $b_{I\backslash i}(C) \cup \{i\}$. Since $\tau$ cannot contain an $I$-broken circuit of $M$, $\tau\backslash i$ does not contain $b_{I\backslash i}(C)$. 

($\supseteq$) Let $\tau$ be a face of $\Delta_w(M\backslash i ,I\backslash i)$ and suppose $C$ is a circuit of $M$. 
If $i\not\in C$, then $C$ is also a circuit of $M\backslash i $, implying that $b_I(C)$ is not contained in $\tau$. 
If $i\in C$ and $C\subseteq I$, then $i = \max(C)$. Since $i$ is not a loop, this implies that $i\in b_I(C)$, which 
 cannot be contained in $\tau$. Similarly, if $i\in C$ and $C\not\subset I$, then $i\in b_I(C)$ and $b_I(C)\not\subset \tau$. 

Finally, let $\tau$ be a face of $\Delta_w(M/ i ,I\backslash i)$ and let $C$ be a circuit of $M$. 
If $i\in C$, then $C\backslash i $ is a circuit of $M/i$. Then 
$b_I(C)$ equals $b_{I\backslash i}(C\backslash i)\cup \{i\}$. Since $\tau$ cannot contain $b_{I\backslash i}(C\backslash i)$, $\tau \cup\{i\}$ does not contain $b_I(C)$. 
If $i\not\in C$, then $C$ is a union of circuits of $M/i$, see \cite[\S~3.1, Exercise~2]{Oxley}. 
If $C\subseteq I$, then there is a circuit $C' \subseteq C$ of $M/i$ containing $\min(C)$. 
Then $C'\subseteq I\backslash i$ and $b_{I\backslash i}(C')$ is a subset of $b_I(C)$. 
Similarly, if $C\not\subseteq I$, then there is a circuit $C' \subseteq C$ of $M/i$ containing $\max(C\backslash I)$, 
giving $b_{I\backslash i}(C') \subseteq b_I(C)$. In either case, $\tau$ is a face of $\Delta_w(M/ i ,I\backslash i)$ and cannot contain
the broken circuit $b_{I\backslash i}(C')$ and therefore $\tau \cup \{i\}$ cannot contain $b_I(C)$.%\qedhere
%\end{enumerate}
\end{proof}



\begin{coro} \label{cor:DeltaDim}
If $M$ is a matroid of rank $d$ with no loops in $I$, then $\Delta_w(M,I)$ is a pure simplicial complex of dimension $d-1$.
\end{coro}

\begin{proof}
We induct on the size of $I$. If $I = \emptyset$, then for every circuit $C$, the broken circuit $b_I(C)$ is 
the maximum element $\max(C)$. In this case, the simplicial complex $\Delta_w(M,I)$ consists of 
one maximal face $B$, where $B$ is the lexicographically smallest basis of $M(\cL)$. 
Here $B$ consists of the elements $i\in [n]$ for which the rank of $[i]$ in $M(\cL)$ is strictly larger than the rank of $[i-1]$. 
Every other element is the maximal element of some circuit of $M$. 

Now suppose $|I| >0$ and consider $i = \max(I)$. If $i$ is a coloop of $M$, then the contraction $M/i$ is a matroid of rank $d-1$ with no loops in $I\backslash i$. 
Then by induction and Theorem~\ref{thm:Complex}\ref{theo3.2_b}~, $\Delta_w(M,I) = \cone (\Delta_w(M/i,I\backslash i),i)$ is a pure simplicial complex of 
dimension $d-1$. Finally, suppose $i$ is neither a loop nor a coloop of $M$. 
Then the deletion $M\backslash i$ is a matroid of rank $d$ and 
no element of $I\backslash i$ is a loop of $M\backslash i$. It follows that $\Delta_w(M\backslash i, I\backslash i)$ is a pure simplicial complex of 
dimension $d-1$. The contraction $M/i$ is a matroid of rank $d-1$, implying that $\Delta_w(M/i,I\backslash i)$ is either empty (if $I\backslash i$ contains a loop of $M/i$), 
or a pure simplicial complex of dimension $d-2$. In either case the decomposition in Theorem~\ref{thm:Complex} finishes the proof.
\end{proof}
The connection between this simplicial complex and the semi-inverted linear space $\inv_I(\cL)$ 
is that when $w \in (\R_+)^n$ has distinct coordinates, the ideal generated by the initial forms $\{ \In_w(f_C) : \text{ $C$ is a circuit of }M(\cL)\}$ 
is the Stanley--Reisner ideal of $\Delta_w(M,I)$. In fact, the initial form of $f_C$ is $\In_w(f_C) =\x^{b_I(C)}$.
The ideal generated by these initial forms is then the Stanley--Reisner ideal $\I_{\Delta_w(M,I) } = \langle \In_w(f_C):C \in \CC (M) \rangle$.


\section{Proof of Theorem~\ref{thm:main}}\label{sec:proof}

In this section, we prove Theorem~\ref{thm:main}. To do this, we first use a flat degeneration of $\inv_I(\cL)$ 
to establish a recursion for its degree. 

\begin{prop}\label{prop:DegreeRecurrence}
Suppose $\cL$ is a linear subspace of $\Cbb^n$ and $I\subseteq [n]$. 
Let $D(\cL,I)$ denote the degree of the affine variety $\inv_I(\cL)$. 
\begin{enumerate}[(\alph*)] 
\item\label{prop4.1_a} %[(a)] 
If $i \in I$ is a loop of $M(\cL)$, then $\inv_I(\cL)$ is empty and $D(\cL,I) = 0$. 
\item\label{prop4.1_b} %[(b)] 
If $i\in I$ is a co-loop of $M(\cL)$, then $D(\cL,I) = D(\cL/i, I\backslash i)$. 
\item\label{prop4.1_c} %[(c)] 
If $i\in I$ is neither a loop nor a coloop of $M(\cL)$ then
\[
D(\cL\backslash i,I\backslash i)  +  D(\cL/i, I\backslash i) \leq D(\cL,I).
\]
\end{enumerate}
\end{prop}
The proof of Theorem~\ref{thm:main} will show that there is actually equality in part~\ref{prop4.1_c}. 

\begin{proof} Without loss of generality, take $i = 1$. 

\ref{prop4.1_a}~If $1\in I$ is a loop of $M(\cL)$ then $\cL$ is contained in the hyperplane $\{x_1=0\}$. 
Therefore the map $\inv_I$ is undefined at every point of $\cL$ and the image $\inv_I(\cL)$ is empty. 
By convention, we take the degree of the empty variety to be zero. 

\ref{prop4.1_b}~If $1$ is a co-loop of $M(\cL)$, then $\cL$ is a direct sum of $\spanrm (e_1)$ and $\cL/1$, meaning that 
any element in $\cL$ can be written as $ae_1 + b$ where $a\in \Cbb$ and $b\in \cL/1$. 
For points at which the map $\inv_I$ is defined, $\inv_I(ae_1 + b) = a^{-1}e_1 + \inv_{I\backslash 1}(b)$. 
From this, we see that $\inv_I(\cL)$ is the direct sum of $\spanrm (e_1)$ and $\inv_{I\backslash 1}(\cL/1)$, 
implying that $\inv_{I}(\cL)$ and $\inv_{I\backslash 1}(\cL/1)$ have the same degree. 

\ref{prop4.1_c}~Let $\I$ denote the ideal of polynomials vanishing on $\inv_I(\cL)$ and 
$\J = \ol{\I}$ denote its homogenization in $\Cbb[x_0,x_1,\ldots , x_n]$. 
Take $w = e_1 \in \R^{n+1}$ and consider $\In_w(\J)$, 
as defined in Section~\ref{subsec:Grobner}. 
We will show that the variety of $\In_w(\J)$ contains 
the image in $\PP^n$ of both $\{0\}\times \inv_{I\backslash 1}(\cL\backslash 1)$ and 
$\mathbb{A}^1(\Cbb)\times \inv_{I\backslash 1}(\cL/ 1)$. 
Since both these varieties have dimension equal to $\dim(\cL)$, the degree of the variety of $\In_w(\J)$ is at least the sum of their degrees. 
The claim then follows by the equality of the Hilbert series of $\J$ and $\In_w(\J)$. 


If $j\in I\backslash 1$ is a loop of $M(\cL)$, then $j$ is a loop of $M(\cL\backslash 1)$ and $D(\cL\backslash 1,I\backslash 1)=0$.
Otherwise the set $U_I$ is Zariski-dense in $\cL$, where $U_I$ denotes the intersection of 
$\cL$ with $(\Cbb^*)^I\times \Cbb^{[n]\backslash I}$. 

Let $\pi_I$ denote the coordinate projection $\Cbb^n \rightarrow \Cbb^I$. On $U_I$, the maps 
$\pi_{I\backslash 1}\circ \inv_I$ and $\inv_{I\backslash 1} \circ \pi_{I\backslash 1}$ are equal:
\[\pi_{I\backslash 1}(\inv_I(x))   =   \inv_{I\backslash 1}(\pi_{I\backslash 1}(x))  = 
\sum_{j\in I\backslash 1} x_j^{-1} e_j + \sum_{j\not\in I} x_j e_j. 
\]
In particular, the points $\inv_{I\backslash 1}(\pi_{I\backslash 1}(U_I))$ are Zariski dense in $\inv_{I\backslash 1}(\cL\backslash 1)$. 
Now let $x$ be a point of $ \inv_I(U_I)$. 
Then $[1: x]$ belongs to the variety of $\J$ and, for every $t\in \Cbb$, the point $(t,t^{e_1}\cdot[1:x])$ belongs to the 
variety of $\ol{\J}^{w}$, as defined in Section~\ref{subsec:Grobner}. Taking $t\rightarrow 0$, we see that $[1:0:\pi_{I\backslash 1}(x)]$ belongs to the variety of $\In_{e_1}(\J)$. 


If $j\in I\backslash 1$ is a loop of $M(\cL/1)$, then $\inv_{I\backslash 1}(\cL/1)$ is empty and 
the claim follows. Otherwise the intersection $U_{I\backslash 1}$ of $\cL/1$ with $\{0\}\times (\Cbb^*)^{I\backslash 1} \times \Cbb^{[n]\backslash I}$ is nonempty and Zariski-dense in $\cL\cap \{x\in \Cbb^n: x_1=0\} \cong \cL/1$. 
Let $x \in U_{I\backslash 1}$. Since $1$ is not a loop of $M(\cL)$, 
there is a point $v \in \cL$ with $v_1 = 1$. Then for any $\lambda, t\in \Cbb^*$,
$x+ (t/\lambda)v$ belongs to $\cL$ and for all but finitely many values of $t$, $y(t) = \inv_I(x+ (t/\lambda)v)$
is defined and has first coordinate $y_1(t) = \lambda/t$. 
Then $[1: y(t)] \in \V(\J)$ and $(t, t^{e_1}\cdot [1:y(t)])$ belongs to $\V\left(\ol{\J}^w\right)$. 
Note that the limit of $t^{e_1}\cdot [1:y(t)] = [1:\lambda : y_2(t): \ldots :y_n(t)]$ as $t\rightarrow 0$ equals $[1:\lambda: \inv_{I\backslash 1}(x)]$. 
Therefore for every point $(\lambda,u)\in \mathbb{A}^1(\Cbb)\times \inv_{I\backslash 1}(\cL/1)$, the point $[1:\lambda: u]$ belongs to $\V\left(\In_{e_1}(\J)\right)$. \end{proof}


We also need the following fact from commutative algebra, included here for completeness, 
which may be clear to readers familiar with the geometry of schemes. 
Recall that for a homogeneous ideal $J \subset \Cbb[\x]$, the Hilbert polynomial of $J$ is 
the polynomial $h(t)$ that agrees with $\dim_\Cbb(\Cbb[\x]/J)_t$ for sufficiently large $t \in \N$. 
Then $ h(t) = \sum_{i=0}^d b_i \binom{t}{d-i}$ for some $d\in \Z_{>0}$, where $b_i\in \Z$ with $b_0>0$.
In a slight abuse of notation, we say that the dimension of $J$ is $\dim(J) = d$ and the degree of $J$ is $\deg(J) = b_0$. 
The ideal $J$ is \emph{equidimensional} of dimension $d$ if $\dim(P)=d$ for every minimal associated prime $P$ of $J$. 



\begin{lemm} \label{lem:EquidimEquality}
Let $I\subseteq J\subseteq\Cbb[\x]$ be equidimensional homogeneous ideals of dimension $d$. If $I$ is radical and $\deg(I)\leq \deg(J)$, then $I$ and $J$ are equal. 
\end{lemm}
\begin{proof}
Let $I = P_1 \cap \cdots \cap P_r$ and $J = Q_1 \cap \cdots \cap Q_s$ be irredundant primary decompositions of $I$ and $J$.
Without loss of generality, we can assume that $\dim(Q_i) = d$ for $1 \leq i \leq u$, and since $\V(J) \subseteq \V(I)$, the prime ideals $P_i$ can be reindexed such that $P_i=\sqrt{Q_i}$, 
implying $Q_i \subseteq P_i$. 
For all $1 \leq i \leq u$, there exists an element $a\in (\cap_{j\neq i}P_j) \cap (\cap_{j\neq i}\sqrt{Q_j})$ with $a\not\in P_i$. 
Then the saturation $I:\langle a\rangle^{\infty} = P_i $ is contained in $ J:\langle a\rangle^{\infty} = Q_i$, implying $P_i = Q_i$. 
This writes the ideal $J$ as $J=P_1\cap \cdots \cap P_{u}\cap Q_{u+1} \cap \cdots \cap Q_s$. 
The degree of an ideal is equal to the sum of the degrees of the top dimensional ideals in its primary decomposition, hence
\[
\deg(I) = \sum_{i=1}^r \deg(P_i)  \quad \text{and} \quad 
\deg(J) = \sum_{i=1}^u \deg(Q_i)=\sum_{i=1}^u \deg(P_i).
\]
The assumption that $\deg(I)\leq \deg(J)$ implies that $r=u$, which gives the reverse containment
$I = P_1 \cap \cdots \cap P_u \supseteq J$.
\end{proof}


\begin{proof}[Proof of Theorem~\ref{thm:main}] 
We proceed by induction on $|I|$. If $|I|=0$, then $\inv_I(\cL)$ is just the linear space $\cL$. 
Then Theorem~\ref{thm:main} reduces to the statement that the linear forms supported on circuits 
form a universal Gr\"obner basis for $\I(\cL)$. See \eg \cite[Prop.~1.6]{GrobnerPolytope}.

Now take $|I|\geq 1$, $w\in (\R_+)^n$ with distinct coordinates, and let $M$ denote the matroid $M(\cL)$.
If $M$ has a loop $i$ in $I$, then for the circuit $C = \{i\}$, the circuit polynomial $f_C$ equals $1$, 
which is a Gr\"obner basis for the ideal of polynomials vanishing on the empty set $\inv_I(\cL)$. 
Therefore we may suppose that $M$ has no loops in $I$, in which case $\inv_I(\cL)$ is a $d$-dimensional 
affine variety of degree $D(\cL,I)$. 

Let $\Delta$ denote the $I$-broken circuit complex $\Delta_w(M,I)$
defined in Section~\ref{sec:SimplicialComplex} and let $\Delta_0$ denote the simplicial complex on elements $\{0,\ldots , n\}$ 
obtained from $\Delta$ by coning over the vertex $0$. Let $\I_{\Delta_0}$ denote the Stanley--Reisner ideal 
of $\Delta_0$, as in Section~\ref{subsec:SRI}. 


Let $\I \subset \Cbb[\x]$ be the ideal of polynomials vanishing on $\inv_I(\cL)$ 
and define the ideal $\J\subset \Cbb[x_0, x_1, \ldots , x_n]$ to be its homogenization with respect to $x_0$. 
Since $\inv_I(\cL)$ is the image of an irreducible variety under a rational map, it is also irreducible. 
It follows that the ideals $\I$ and $\J$ are prime. 
For a circuit polynomial $f_C$, its homogenization $\ol{f_C}$ belongs to $\J$ and since $w\in (\R_+)^n$, 
\[
\In_{(0,w)}(\ol{f_C}) = \In_{w}(f_C) = \begin{cases}
 a_k \x^{C\backslash k} & \text{ if $C \subseteq I $ and $k = \argmin \{w_j: j\in C\}$ } \\ 
 a_k \x^{C\cap I\cup k} & \text{ if $C \not\subseteq I $ and $k = \argmax \{w_j: j\in C\backslash I\}$. } \\ 
\end{cases}
\]
Up to a scalar multiple, $\In_{w}(f_C)$ equals the 
square-free monomial corresponding to the $I$-broken circuit of $C$, namely $\x^{b_I(C)}$. 
It follows that 
\[
\langle \In_{w}(f_C) \reldp C \in \CC (M) \rangle = \I_{\Delta}
\ \ 
\text{ and } 
\ \ 
\langle \In_{(0,w)}(\ol{f_C}) \reldp C \in \CC (M) \rangle = \I_{\Delta_0}. 
\]
From this we see that $\I_{\Delta_0} \subseteq \In_{(0,w)}(\J)$. 


Let $i = \argmax \{w_j : j\in I\}$. By the inductive hypothesis, $D(\cL\backslash i, I\backslash i)$ and $D(\cL/ i, I\backslash i)$
are the number of facets of $\Delta_w(M\backslash i, I\backslash i)$ and $\Delta_w(M/i, I\backslash i)$, respectively. 
Therefore by Theorem~\ref{thm:Complex}, $\Delta$ and thus $\Delta_0$ each have 
$D(\cL/ i, I\backslash i)$ facets if $i$ is a coloop of $M$ and $D(\cL\backslash i, I\backslash i) + D(\cL/ i, I\backslash i)$ facets otherwise. 
Then by Proposition~\ref{prop:DegreeRecurrence}, $\Delta_0$ has at most $D(\cL, I)$ facets and the Stanley--Reisner ideal 
$\I_{\Delta_0}$ has degree $\leq D(\cL, I)$. 


Since $\Delta_0$ is a pure simplicial complex of dimension $d$, $\I_{\Delta_0}$ is an 
equidimensional ideal of dimension $d$. 
As $\J$ is a prime $d$-dimensional ideal, its initial ideal $\In_{(0,w)}(\J)$ is equidimensional of the same 
dimension, see \cite[Lemma~2.4.12]{tropBook}. 

The ideals $\I_{\Delta_0}$ and $\In_{(0,w)}(\J)$ then satisfy the hypotheses of Lemma~\ref{lem:EquidimEquality}, 
and we conclude that they are equal. By \cite[Prop.~2.6.1]{tropBook}, restricting to $x_0=1$ 
gives that 
\[
\I_{\Delta}  =  \langle \In_w(f_C) \reldp  C\in \CC (M) \rangle =  \In_w(\I). 
\]
As this holds for every $w\in (\R_+)^n$ with distinct coordinates, it will also hold for arbitrary $w \in (\R_{\geq 0})^n$ (see \cite[Prop.~1.13]{GrobnerPolytope}).
It follows from \cite[Cor.~1.9, 1.10]{GrobnerPolytope} that the circuit polynomials $\{f_C : C\in \CC (M)\}$ form a universal Gr\"obner basis for $\I$. 
\end{proof}
\begin{thebibliography}{99}
\bibitem[1]{AB16} Ardila, Federico; Boocher, Adam The closure of a linear space in a product of lines, J. Algebraic Combin., Volume 43 (2016) no. 1, pp. 199-235.
\bibitem[5]{FSW17} Fink, Alex; Speyer, David E.; Woo, Alexander A Gröbner basis for the graph of the reciprocal plane (2017) (https://arxiv.org/abs/1703.05967).
\bibitem[8]{tropBook} Maclagan, Diane; Sturmfels, Bernd Introduction to tropical geometry, Graduate Studies in Mathematics, 161, American Mathematical Society, Providence, RI, 2015, xii+363 pages.
\bibitem[11]{Oxley} James Oxley, Matroid theory, second ed., Oxford Graduate Texts in Mathematics, vol.21, Oxford University Press, Oxford, 2011.
\bibitem[12]{PS} Proudfoot, Nicholas; Speyer, David A broken circuit ring, Beiträge Algebra Geom., Volume 47 (2006) no. 1, pp. 161-166.
\bibitem[16]{StanleyCCABook} Richard P. Stanley, Combinatorics and commutative algebra, second ed., Progress in Mathematics, vol.41, Birkhäuser Boston, Inc., Boston, MA, 1996.
\bibitem[17]{GrobnerPolytope} Sturmfels, Bernd Gröbner bases and convex polytopes, University Lecture Series, 8, American Mathematical Society, Providence, RI, 1996, xii+162 pages.
\end{thebibliography}
\end{document}
